% ============================================================
%  Chương V – Triển khai và ứng dụng thực tế
%  ~1200 từ · 3.5 trang
% ============================================================
\chapter{Triển khai và ứng dụng thực tế}
\label{chap:deploy}

% Ý chính: PKI không tồn tại trong lý thuyết — đây là hạ tầng vô hình
% mà mọi giao dịch số đang phụ thuộc vào từng giây.

% ----------------------------------------------------------
\section{HTTPS / TLS}
\label{sec:5.1}

\subsubsection*{Lịch sử SSL đến TLS 1.3}
% Ý chính: SSL 2.0 (1995, broken) → SSL 3.0 (POODLE) → TLS 1.0/1.1 (deprecated)
% → TLS 1.2 (2008, vẫn phổ biến) → TLS 1.3 (2018 \cite{RFC8446}, cải tiến lớn)

\subsubsection*{TLS Handshake và xác minh chứng chỉ}
% Ý chính: ClientHello → ServerHello + Certificate → CertificateVerify → Finished
% 4 bước browser validation: chain lên root, validity period, revocation, SAN match domain

\subsubsection*{Cải tiến trong TLS 1.3}
% Ý chính: 1-RTT (giảm latency), PFS bắt buộc (only ECDHE/DHE),
% loại bỏ hoàn toàn RSA key exchange, RC4, 3DES, SHA-1 trong handshake

% ----------------------------------------------------------
\section{Chữ ký số văn bản điện tử}
\label{sec:5.2}

\subsubsection*{Nguyên lý ký số văn bản}
% Ý chính: Hash(document) + Sign(private key) → chữ ký số đính kèm tài liệu
% Phân biệt ký số (cryptographic) khỏi chữ ký điện tử (scanned image)

\subsubsection*{Khung pháp lý tại Việt Nam}
% Ý chính: Luật Giao dịch điện tử 2023, Nghị định 130/2018
% Ứng dụng: hóa đơn điện tử (Nghị định 123/2020), khai báo hải quan, văn bản hành chính

\subsubsection*{CA nhà nước Việt Nam}
% Ý chính: NEAC (Root CA quốc gia dưới Bộ TT\&TT)
% Các CA cấp 2: VNPT-CA, Viettel-CA, BKAV-CA, FPT-CA

% ----------------------------------------------------------
\section{S/MIME -- Email bảo mật}
\label{sec:5.3}
% Ý chính: ký số và mã hóa email theo chuẩn \cite{RFC8551}
% Hybrid encryption: mã hoá body bằng AES (session key), mã hoá session key bằng RSA
% Tích hợp trong Outlook, Thunderbird; cert cần từ CA được trust store tin

% ----------------------------------------------------------
\section{VPN và IPSec}
\label{sec:5.4}
% Ý chính: IKEv2 dùng cert X.509 để xác thực VPN gateway và client
% Ưu điểm so với PSK: không còn shared secret để lộ; mỗi user có cert riêng dễ revoke
% OpenVPN với full PKI: Root CA → Client cert + Server cert

% ----------------------------------------------------------
\section{Code Signing}
\label{sec:5.5}
% Ý chính: bảo vệ chuỗi phân phối phần mềm — đảm bảo binary không bị tamper
% Windows Authenticode (Portable Executable), macOS Gatekeeper/Notarization, Android APK
% EV Code Signing: được CA xác minh kỹ hơn → bỏ qua SmartScreen ngay lập tức

% ----------------------------------------------------------
\section{PKI trong eGovernment Việt Nam}
\label{sec:5.6}
% Ý chính: ứng dụng PKI ở quy mô quốc gia
% Cổng Dịch vụ công quốc gia (DVCQG), một cửa điện tử, eTax, e-Invoice
% Kiến trúc: NEAC Root CA → CA Bộ/Tỉnh → CA Doanh nghiệp/Cá nhân
% Thực trạng: tốc độ triển khai đang tăng nhanh nhưng vẫn còn thách thức interoperability
